% ══════════════════════════════════════════════════════════════
%  Section 9.5: Transaction Lifecycle --- Create → Sign → Verify → Mempool
% ══════════════════════════════════════════════════════════════
\chapter{Chain: Transaction Lifecycle}
\label{ch:transaction-lifecycle}

\begin{learnbox}
\begin{itemize}
  \item How a transaction flows from wallet creation through network propagation to mining boundary
  \item Transaction construction: selecting inputs from UTXO set and creating outputs
  \item Signing: the ``trimmed copy'' pattern that authorizes spending without circular signatures
  \item Verification: checking signatures at the mining boundary
  \item Mempool admission: accepting pending transactions and avoiding local conflicts
  \item Mining trigger: when and how the system decides to start mining
\end{itemize}
\end{learnbox}

When you want to send coins to someone, what actually happens behind the scenes? This section traces a \index{transaction}transaction through its complete \index{transaction!lifecycle}lifecycle---from the moment you decide to send coins, through \index{construction}construction and \index{digital signature!signing}signing, \index{verification}verification and \index{mempool}mempool admission, to \index{network!propagation}network propagation and eventual inclusion in a \index{block}block.

% ──────────────────────────────────────────────────────────────
\section{Overview: The Complete Transaction Journey}
\label{sec:ch09-5-overview}

A \index{transaction}transaction goes through several distinct \index{stage}stages, each handled by different parts of the \index{system}system:

\begin{figure}[htbp]
\centering
\begin{tikzpicture}[
  node distance = 1.0cm,
  stage/.style = {rectangle, rounded corners=4pt, draw=sectioncolor, fill=sectioncolor!10,
                  text width=4.2cm, align=center, minimum height=1.1cm, font=\sffamily\small},
  code/.style = {font=\ttfamily\footnotesize, color=codekw},
  arrow/.style = {-{Stealth[length=4pt,width=3pt]}, thick, color=brand},
]

% Wallet stage
\node[stage] (wallet) at (0, 6) {
  \textbf{Wallet} \\
  \small construct \& sign
};
\node[font=\tiny\sffamily, right=0.3cm of wallet, text width=2cm, align=left] {
  \code{new\_utxo\_tx}\\
  \code{find\_spend}\\
  \code{sign}
};

% Node mempool stage
\node[stage] (mempool) at (0, 4.5) {
  \textbf{Node} \\
  \small add to mempool
};
\node[font=\tiny\sffamily, right=0.3cm of mempool, text width=2cm, align=left] {
  \code{process\_tx}\\
  \code{add\_to\_memory}
};

% Network broadcast stage
\node[stage] (network) at (0, 3) {
  \textbf{Network} \\
  \small broadcast peers
};
\node[font=\tiny\sffamily, right=0.3cm of network, text width=2cm, align=left] {
  \code{send\_inv}\\
  (Tx, txid)
};

% Mining verification stage
\node[stage] (mining) at (0, 1.5) {
  \textbf{Mining} \\
  \small verify transaction
};
\node[font=\tiny\sffamily, right=0.3cm of mining, text width=2cm, align=left] {
  \code{mine\_block}\\
  \code{verify}
};

% Confirmation stage
\node[stage] (confirm) at (0, 0) {
  \textbf{Confirmation} \\
  \small in blockchain
};
\node[font=\tiny\sffamily, right=0.3cm of confirm, text width=2cm, align=left] {
  \code{remove\_from}\\
  \code{memory\_pool}
};

% Arrows
\draw[arrow] (wallet) -- (mempool);
\draw[arrow] (mempool) -- (network);
\draw[arrow] (network) -- (mining);
\draw[arrow] (mining) -- (confirm);

\end{tikzpicture}
\caption{Transaction Lifecycle: From Wallet to Blockchain Confirmation}
\label{fig:ch09-5-tx-lifecycle}
\end{figure}

The complete flow: \textbf{\index{construction}construct \(\to\) \index{digital signature!signing}sign \(\to\) \index{submission}submit \(\to\) \index{mempool}mempool \(\to\) (async) \index{broadcast}broadcast/\index{mining}mining \(\to\) \index{verification}verify at \index{mining boundary}mining boundary \(\to\) \index{confirmation}confirmation}.

% ──────────────────────────────────────────────────────────────
\section{Step 1: Entry Point --- Top-Level Send Payment}
\label{sec:ch09-5-step1}

This is the highest-level \index{orchestration}orchestration method: it creates a \index{transaction}transaction and hands it to the \index{node}node for \index{mempool}mempool/\index{networking}network processing.

\begin{listing}[htbp]
\caption{\code{NodeContext::btc\_transaction} --- send payment entry point}
\label{lst:ch09-5-2}
\begin{minted}[fontsize=\footnotesize]{rust}
// Source: context.rs
pub async fn btc_transaction(
    &self,
    wlt_frm_addr: &WalletAddress,
    wlt_to_addr: &WalletAddress,
    amount: i32,
) -> Result<String> {
    // Create a UTXOSet handle (derived-state accessor).
    // Why: transaction construction needs coin selection;
    // signing needs previous-tx lookups via chainstate.
    let utxo_set = UTXOSet::new(self.blockchain.clone());

    // Build + sign the transaction (wallet-side responsibility).
    // Why: the node should not ``invent'' spends;
    // it only accepts/relays what the wallet constructed.
    let utxo = Transaction::new_utxo_transaction(
        wlt_frm_addr, wlt_to_addr, amount, &utxo_set,
    ).await?;

    // Submit to the node's mempool/network pipeline.
    // Why: once constructed, a tx must enter the node's
    // pending set before it can be mined/relayed.
    let addr_from = crate::GLOBAL_CONFIG.get_node_addr();
    self.process_transaction(&addr_from, utxo).await
}
\end{minted}
\end{listing}

\begin{notebox}
This method separates responsibilities: wallet-side work (construction + signing) is followed by node-side work (mempool/network processing). The return value is the transaction ID for the caller to track.
\end{notebox}

% ──────────────────────────────────────────────────────────────
\section{Step 2: Construct a Spending Transaction}
\label{sec:ch09-5-step2}

Now we see how the \index{transaction}transaction is actually built. This is where the \index{wallet}wallet selects which coins to spend, creates the payment and \index{change output}change outputs, computes the \index{txid}transaction ID, and \index{digital signature!signing}signs everything.

\begin{listing}[htbp]
\caption{\code{Transaction::new\_utxo\_transaction} --- build and sign}
\label{lst:ch09-5-3}
\begin{minted}[fontsize=\footnotesize]{rust}
// Source: transaction.rs
pub async fn new_utxo_transaction(
    from_wlt_addr: &WalletAddress,
    to_wlt_addr: &WalletAddress,
    tx_amount: i32,
    utxo_set: &UTXOSet,
) -> Result<Transaction> {
    let wallets = WalletService::new()?;
    let from_wallet = wallets.get_wallet(from_wlt_addr)?;
    let from_pub_key_hash = hash_pub_key(
        from_wallet.get_public_key()
    );
    let (available_funds, valid_outputs) = utxo_set
        .find_spendable_outputs(&from_pub_key_hash, tx_amount)
        .await?;

    if available_funds < tx_amount {
        return Err(BtcError::NotEnoughFunds);
    }

    let mut inputs = vec![];
    for (txid_hex, out_indexes) in valid_outputs {
        let txid = HEXLOWER.decode(txid_hex.as_bytes())?;
        for vout_idx in out_indexes {
            inputs.push(TXInput::new(&txid, vout_idx));
        }
    }

    let mut outputs = vec![TXOutput::new(tx_amount, to_wlt_addr)?];
    if available_funds > tx_amount {
        let change = available_funds - tx_amount;
        outputs.push(TXOutput::new(change, from_wlt_addr)?);
    }

    let mut tx = Transaction {
        id: vec![],
        vin: inputs,
        vout: outputs,
    };
    tx.id = tx.hash()?;
    tx.sign(utxo_set.get_blockchain(), from_wallet.get_pkcs8())
        .await?;
    Ok(tx)
}
\end{minted}
\end{listing}

\begin{checkpointbox}
This method:
\begin{enumerate}
  \item Loads the sender's \index{key}keys and derives the sender's \code{\index{public key}pub\_key\_hash}
  \item Performs \index{coin selection}coin selection via \code{\index{UTXOSet!find\_spendable\_outputs}find\_spendable\_outputs}
  \item Turns selected \index{outpoint}outpoints into \code{\index{TXInput}TXInput}s (empty \index{digital signature}signatures for now)
  \item Builds \code{\index{TXOutput}TXOutput}s: payment plus optional \index{change output}change
  \item Finalizes by computing \index{txid}txid (\code{\index{hash}tx.hash()}) and \index{digital signature!signing}signing (\code{\index{sign}tx.sign(...)})
\end{enumerate}

Change is a \emph{\index{new output}new output}, not ``leftover in an account'', because \index{UTXO}UTXO-based systems must explicitly create all spendable coins.
\end{checkpointbox}

% ──────────────────────────────────────────────────────────────
\section{Step 3: Sign and Verify (Trimmed Copy Pattern)}
\label{sec:ch09-5-step3}
\index{digital signature!signing}\index{trimmed copy}

Once the \index{transaction}transaction is constructed, it must be \index{authorization}authorized. \index{digital signature}Signatures prove that the sender owns the coins being spent.

\begin{listing}[htbp]
\caption{Trimmed copy, signing, and verification}
\label{lst:ch09-5-4}
\begin{minted}[fontsize=\footnotesize]{rust}
// Source: transaction.rs
fn trimmed_copy(&self) -> Transaction {
    let inputs = self.vin.iter()
        .map(|i| TXInput::new(i.get_txid(), i.get_vout()))
        .collect();
    Transaction {
        id: self.id.clone(),
        vin: inputs,
        vout: self.vout.clone()
    }
}

async fn sign(
    &mut self,
    blockchain: &BlockchainService,
    private_key: &[u8],
) -> Result<()> {
    let mut tx_copy = self.trimmed_copy();
    for (idx, vin) in self.vin.iter_mut().enumerate() {
        let prev_tx = blockchain
            .find_transaction(vin.get_txid())
            .await??;

        // Temporarily inject the referenced output's lock
        tx_copy.vin[idx].pub_key =
            prev_tx.vout[vin.vout].pub_key_hash.clone();
        tx_copy.id = tx_copy.hash()?;
        let signature = schnorr_sign_digest(
            private_key,
            tx_copy.get_id()
        )?;
        vin.signature = signature;
        tx_copy.vin[idx].pub_key = vec![];
    }
    Ok(())
}

pub async fn verify(
    &self,
    blockchain: &BlockchainService
) -> Result<bool> {
    if self.is_coinbase() { return Ok(true); }
    let mut trimmed_copy = self.trimmed_copy();
    for (idx, vin) in self.vin.iter().enumerate() {
        let prev_tx = blockchain
            .find_transaction(vin.get_txid())
            .await??;
        let pkh = prev_tx.vout[vin.vout]
            .pub_key_hash.clone();
        trimmed_copy.vin[idx].pub_key = pkh;
        trimmed_copy.id = trimmed_copy.hash()?;
        let is_valid = schnorr_sign_verify(
            vin.get_pub_key(),
            vin.get_signature(),
            trimmed_copy.get_id(),
        );
        if !is_valid {
            return Ok(false);
        }
        trimmed_copy.vin[idx].pub_key = vec![];
    }
    Ok(true)
}
\end{minted}
\end{listing}

\begin{notebox}
\textbf{The trimmed copy pattern} avoids circular signing (the signature cannot sign itself):
\begin{enumerate}
  \item Create a copy with signatures cleared and other signing-related fields blank
  \item For each input: inject the referenced output's \code{pub\_key\_hash} to bind the signature to the exact outpoint being spent
  \item Hash that temporary copy
  \item Sign the hash (which represents the transaction \emph{shape} and referenced outputs)
  \item Verification recomputes that same digest and verifies the signature
\end{enumerate}

This ensures signatures commit to the exact outpoint and locking data, preventing signature reuse across different spends.
\end{notebox}

% ──────────────────────────────────────────────────────────────
\section{Step 4: Accept into Mempool and Spawn Background Work}
\label{sec:ch09-5-step4}

This is the node's acceptance boundary: reject duplicates, add to the mempool, then spawn propagation/mining trigger in the background.

\begin{listing}[htbp]
\caption{\code{NodeContext::process\_transaction} --- mempool acceptance}
\label{lst:ch09-5-5}
\begin{minted}[fontsize=\footnotesize]{rust}
// Source: context.rs
pub async fn process_transaction(
    &self,
    addr_from: &std::net::SocketAddr,
    utxo: Transaction,
) -> Result<String> {
    if transaction_exists_in_pool(&utxo) {
        let txid = utxo.get_tx_id_hex();
        return Err(
            BtcError::TransactionAlreadyExistsInMemoryPool(
                txid
            )
        );
    }
    add_to_memory_pool(utxo.clone(), &self.blockchain)
        .await?;

    let context = self.clone();
    let addr_copy = *addr_from;
    let tx = utxo.clone();
    tokio::spawn(async move {
        let _ = context
            .submit_transaction_for_mining(&addr_copy, tx)
            .await;
    });
    Ok(utxo.get_tx_id_hex())
}
\end{minted}
\end{listing}

\begin{notebox}
This method:
\begin{itemize}
  \item \textbf{Synchronously} checks for duplicate (by \code{txid})
  \item \textbf{Synchronously} adds to mempool via \code{add\_to\_memory\_pool}
  \item \textbf{Asynchronously} (spawns background task) propagates and triggers mining
  \item Returns the \code{txid} immediately to the caller
\end{itemize}

Important: this method does \textbf{not} call \code{Transaction::verify} before storing. Signature verification is enforced at the \textbf{mining boundary} (\code{BlockchainService::mine\_block}). This is a learning simplification; production systems validate at mempool admission too.
\end{notebox}

% ──────────────────────────────────────────────────────────────
\section{Step 5: Mempool Internals}
\label{sec:ch09-5-step5}

How does the mempool store, track, and manage pending transactions?

\begin{listing}[htbp]
\caption{Mempool operations: add, remove, check existence}
\label{lst:ch09-5-6}
\begin{minted}[fontsize=\footnotesize]{rust}
// Source: txmempool.rs
pub async fn add_to_memory_pool(
    tx: Transaction,
    blockchain: &BlockchainService,
) -> Result<()> {
    let tx_id = tx.get_tx_id_hex();
    if GLOBAL_MEMORY_POOL.lock().await.contains_key(&tx_id) {
        return Err(BtcError::TransactionAlreadyExistsInMemoryPool(
            tx_id
        ));
    }
    GLOBAL_MEMORY_POOL
        .lock()
        .await
        .insert(tx_id, tx.clone());

    // Mark referenced inputs as reserved
    let utxo_set = UTXOSet::new(blockchain.clone());
    for input in tx.get_vin() {
        if let Ok(Some(mut outputs)) =
            utxo_set.find_utxo_by_txid(input.get_txid()).await
        {
            if input.vout < outputs.len() {
                outputs[input.vout]
                    .set_in_global_mem_pool(true);
            }
        }
    }
    Ok(())
}

pub fn transaction_exists_in_pool(tx: &Transaction) -> bool {
    let tx_id = tx.get_tx_id_hex();
    GLOBAL_MEMORY_POOL.blocking_lock().contains_key(&tx_id)
}

pub async fn remove_from_memory_pool(
    tx: &Transaction,
    blockchain: &BlockchainService,
) -> Result<()> {
    let tx_id = tx.get_tx_id_hex();
    GLOBAL_MEMORY_POOL.lock().await.remove(&tx_id);

    // Unmark referenced inputs
    let utxo_set = UTXOSet::new(blockchain.clone());
    for input in tx.get_vin() {
        if let Ok(Some(mut outputs)) =
            utxo_set.find_utxo_by_txid(input.get_txid()).await
        {
            if input.vout < outputs.len() {
                outputs[input.vout]
                    .set_in_global_mem_pool(false);
            }
        }
    }
    Ok(())
}
\end{minted}
\end{listing}

\begin{notebox}
The mempool is a global in-memory map of \code{txid -> Transaction}. When a transaction is added:
\begin{itemize}
  \item It is stored in \code{GLOBAL\_MEMORY\_POOL}
  \item Its referenced inputs are marked with \code{in\_global\_mem\_pool = true} to reserve them locally
\end{itemize}

When removed (after confirmation):
\begin{itemize}
  \item It is deleted from the mempool
  \item Its referenced inputs are unmarked
\end{itemize}

This reservation mechanism reduces (but does not eliminate) the chance that the node builds conflicting pending spends.
\end{notebox}

% ──────────────────────────────────────────────────────────────
\section{Step 6: Propagation and Mining Trigger}
\label{sec:ch09-5-step6}

After mempool admission, a background task propagates the transaction to peers and may trigger mining.

\begin{listing}[htbp]
\caption{Background task: submit for mining and relay}
\label{lst:ch09-5-7}
\begin{minted}[fontsize=\footnotesize]{rust}
// Source: context.rs (simplified)
async fn submit_transaction_for_mining(
    &self,
    addr_from: &std::net::SocketAddr,
    tx: Transaction,
) -> Result<()> {
    // Broadcast transaction to peers by sending INV message
    self.send_inv_for_transaction(&tx).await?;

    // If this is a mining node and mempool
    // is large enough, trigger mining
    let mempool_size = self.get_mempool_size()?;
    if self.is_mining_enabled()
        && mempool_size >= MINING_TRIGGER_THRESHOLD
    {
        self.trigger_mining_with_candidate_txs().await?;
    }

    Ok(())
}
\end{minted}
\end{listing}

\begin{notebox}
The background task handles two concerns:
\begin{enumerate}
  \item \textbf{Propagation}: announces the transaction to peers via INV (inventory) message
  \item \textbf{Mining trigger}: if mining is enabled and mempool exceeds a threshold, initiates mining with the pending transactions
\end{enumerate}

This decoupling (background task) allows the initial mempool submission to return immediately without waiting for network I/O or mining computation.
\end{notebox}

% ──────────────────────────────────────────────────────────────
\needspace{3in}
\section{Step 7: Mining Boundary --- Defensive Verification}
\label{sec:ch09-5-step7}

Before a transaction can be included in a block, it is verified again at the mining boundary.

\begin{listing}[htbp]
\caption{Mining boundary: verify each transaction before mining}
\label{lst:ch09-5-8}
\begin{minted}[fontsize=\footnotesize]{rust}
// Source: chainstate.rs
pub async fn mine_block(
    &self,
    transactions: &[Transaction],
) -> Result<Block> {
    // Gate: validate before mutating state.
    for transaction in transactions {
        let is_valid = transaction.verify(self).await?;
        if !is_valid {
            return Err(BtcError::InvalidTransaction);
        }
    }

    // Acquire write lock
    let blockchain_guard = self.0.write().await;

    // Re-validate transaction inputs under the write lock
    // (stale mining protection---see Chapter 11)
    let db = blockchain_guard.get_db();
    let utxo_tree = db.open_tree("chainstate")?;
    for tx in transactions {
        if tx.is_coinbase() { continue; }
        for input in tx.get_vin() {
            // Verify each input's UTXO still exists...
        }
    }

    // Mining + persistence happens inside the storage layer.
    blockchain_guard.mine_block(transactions).await
}
\end{minted}
\end{listing}

\begin{notebox}
\textbf{The mining boundary} is the final defensive gate before blocks are created:
\begin{enumerate}
  \item \textbf{Pre-lock verification}: each transaction signature is checked before acquiring the exclusive write lock
  \item \textbf{Lock acquisition}: the write lock is taken to ensure atomic block creation
  \item \textbf{Stale-mining protection}: transaction inputs are re-checked under the lock to catch any inputs that were spent by a competing block between pre-lock and lock acquisition
  \item \textbf{Persistence}: only after all checks pass is the block actually persisted and UTXO state updated
\end{enumerate}

This double-check pattern prevents the node from creating blocks with invalid or already-spent transactions.
\end{notebox}

% ──────────────────────────────────────────────────────────────
\section{Step 8: Mining Trigger Entry Points}
\label{sec:ch09-5-step8}

The system decides to start mining in response to two events:

\begin{listing}[htbp]
\caption{Mining trigger 1: mempool threshold}
\label{lst:ch09-5-9}
\begin{minted}[fontsize=\footnotesize]{rust}
// Trigger 1: Mempool size reaches threshold
if is_mining_enabled()
    && get_mempool_size()
        >= TRANSACTION_THRESHOLD
{
    prepare_candidate_txs_and_mine().await?;
}
\end{minted}
\end{listing}

\begin{listing}[htbp]
\caption{Mining trigger 2: on-demand mining (manual trigger)}
\label{lst:ch09-5-10}
\begin{minted}[fontsize=\footnotesize]{rust}
// Trigger 2: Manual/on-demand mining
// (e.g., via API: POST /api/v1/mining/generatetoaddress)
pub async fn generate_to_address(
    State(node): State<Arc<NodeContext>>,
    Json(req): Json<GenerateToAddressRequest>,
) -> Result<Json<Block>> {
    let block = node
        .mine_block_to_address(&req.address)
        .await?;
    Ok(Json(block))
}
\end{minted}
\end{listing}

\begin{notebox}
Mining is triggered in two ways:
\begin{enumerate}
  \item \textbf{Automatic}: when the mempool reaches a size threshold (e.g., during transaction flooding)
  \item \textbf{On-demand}: via API request or other explicit signals (e.g., ``mine now to address X'')
\end{enumerate}

Once mining is triggered, the system prepares candidate transactions from the mempool, verifies them, and calls \code{BlockchainService::mine\_block}, which performs the final verification and persistent storage.
\end{notebox}

% ──────────────────────────────────────────────────────────────
\section{Summary}
\label{sec:ch09-5-summary}

\begin{chaptersummary}
\item \textbf{Entry point}: \code{NodeContext::btc\_transaction} orchestrates wallet-side (construction + signing) and node-side (mempool/network) work
\item \textbf{Construction}: \code{Transaction::new\_utxo\_transaction} selects inputs from UTXO set and creates outputs (payment + optional change)
\item \textbf{Signing}: the \textbf{trimmed copy pattern} computes signatures without circular self-reference by temporarily injecting locking data, hashing, and signing
\item \textbf{Mempool admission}: \code{process\_transaction} checks for duplicates, stores the transaction, and spawns background propagation/mining trigger
\item \textbf{Reservation}: the \code{in\_global\_mem\_pool} flag is set on referenced outputs to reduce local double-selection
\item \textbf{Propagation}: asynchronous task sends INV messages to peers to announce the transaction
\item \textbf{Mining trigger}: automatic (mempool threshold) or on-demand (API request) initiation of block mining
\item \textbf{Mining boundary}: \code{BlockchainService::mine\_block} performs defensive verification before creating and persisting the block
\item \textbf{Stale-mining protection}: transaction inputs are re-verified under the write lock to catch competing blocks
\item \textbf{Confirmation}: once a block is accepted, its transactions are removed from the mempool
\end{chaptersummary}